\clearpage
\phantomsection
\chapter{IMPLEMENTATION}

\section{Development Environment}

\begin{table}[H]
    \centering
    \caption{Development environment specifications.}
    \label{tab:dev_env}
    \begin{tabular}{ll}
        \toprule
        \textbf{Component} & \textbf{Version / Specification} \\
        \midrule
        Game Engine Language  & C\# 13.0 \\
        Runtime Framework     & .NET 9.0 (cross-platform) \\
        Python Version        & 3.10 \\
        RL Framework          & Stable-Baselines3 2.3.0 \\
        RL Extension          & sb3-contrib 2.3.0 (MaskablePPO) \\
        Deep Learning Backend & PyTorch 2.2.0 + CUDA 12.1 \\
        Python-C\# Bridge     & pythonnet 3.0.3 \\
        IDE                   & Visual Studio Code 1.87 \\
        Version Control       & Git / GitHub \\
        Training Platform     & Kaggle Notebooks (GPU T4 $\times$2) \\
        Local OS              & Windows 11 (development) \\
        Training OS           & Linux (Kaggle, Ubuntu 22.04) \\
        \bottomrule
    \end{tabular}
\end{table}

The choice of C\# as the game engine language was motivated by its strong type system, high simulation throughput (the pure C\# game loop executes $\sim$2,500 complete runs per second on a single core), and .NET 9's native cross-platform support. Python was selected for the RL and optimization layer due to the maturity of the PyTorch and Stable-Baselines3 ecosystems.

\section{Core Game Engine Implementation}

\subsection{Project Structure}

\begin{lstlisting}[language=bash, caption={C\# project structure}]
src/
|-- Roguelike.Core/                 # Shared game library (.dll)
|   |-- AI/
|   |   |-- HeuristicPlayerAI.cs   # Deterministic heuristic agent
|   |   |-- IPlayerAgent.cs        # Interface for all agents
|   |-- Combat/
|   |   |-- CombatManager.cs       # Turn-by-turn combat logic
|   |   |-- DeckManager.cs         # Draw/discard pile management
|   |   |-- ActiveEffect.cs        # Status effect runtime
|   |-- Map/
|   |   |-- MapGenerator.cs        # Procedural map generation
|   |   |-- RoomHandlers/          # One handler per room type
|   |-- Data/
|   |   |-- CardData.cs, EnemyData.cs, RelicData.cs, ...
|   |   |-- GameDataLoader.cs      # JSON deserialization
|   |-- Optimization/
|   |   |-- GeneticAlgorithm.cs
|   |   |-- NSGA2Optimizer.cs
|   |   |-- BalanceGenome.cs
|   |   |-- HierarchicalGenome.cs
|   |   |-- FitnessEvaluator.cs
|   |-- GameController.cs          # Top-level game run orchestrator
|-- Roguelike/                     # Console entry point
|   |-- Program.cs
|-- Roguelike.Tests/               # xUnit test suite
GameData.json                      # All game entities
\end{lstlisting}

\subsection{Game Controller}

\texttt{GameController.cs} is the top-level orchestrator. It manages the game state machine, transitioning between the following states: \texttt{MapNavigation}, \texttt{Combat}, \texttt{CardReward}, \texttt{RelicReward}, \texttt{ShopVisit}, \texttt{EventRoom}, and \texttt{RestSite}. It owns a reference to the active \texttt{IPlayerAgent} and calls \texttt{GetDecision(GameRun)} at each decision point.

The \texttt{GameRun} object is a snapshot of the complete game state passed to the agent. It is a value-type aggregate (not a reference to the mutable game state) to prevent agents from modifying the game state directly.

\subsection{Combat System}

\texttt{CombatManager.ExecuteTurn()} implements the following steps each turn:

\begin{enumerate}
    \item Reset mana to base value; draw 5 cards into hand via \texttt{DeckManager.DrawCards()}.
    \item Apply start-of-turn relic effects (e.g., ``draw 2 extra cards'').
    \item Call \texttt{IPlayerAgent.GetDecision()} in a loop until the agent returns \texttt{EndTurn}.
    \item For each \texttt{PlayCard} decision: validate against action mask, deduct mana, execute card actions (damage, block, effects), move card to discard pile.
    \item Execute each living enemy's scheduled action via \texttt{ExecuteEnemyAction()}.
    \item Apply damage: reduce block first, then HP; trigger thorns and on-damage relics.
    \item Check win/loss conditions; if all enemies defeated, return \texttt{CombatResult.Victory}.
\end{enumerate}

Status effects (Strength, Weakness, Poison, Burn, Vulnerability) are managed by \texttt{ActiveEffect.cs}. Each effect has an intensity and duration; they are applied during the damage calculation step and decremented at end-of-turn.

\subsection{Map Generation}

The map consists of a directed acyclic graph of 15 floors, with 2--3 branches per floor. Map generation proceeds as follows:

\begin{enumerate}
    \item Start with Floor 15 as the fixed ``Boss'' room.
    \item For each floor from 14 down to 1, generate 2--3 nodes.
    \item Assign room types to each node by sampling from a categorical distribution parameterized by \texttt{RoomTypeWeights} (part of the balance genome).
    \item Apply structural constraints: at least one Rest room in each Act; no two consecutive Elite rooms; no Shop on Floor 1.
    \item Connect nodes to ensure full connectivity from Floor 1 to Floor 15.
\end{enumerate}

\subsection{Data-Driven Design}

All game entities are defined in \texttt{GameData.json} and loaded at startup via \texttt{GameDataLoader.cs} using System.Text.Json. The \texttt{GenomeApplicator} class takes a genome and produces a modified copy of \texttt{GameData} without mutating the original, ensuring thread safety when multiple simulations run in parallel.

An example card definition is shown below:

\begin{lstlisting}[language=C, caption={Example card definition in GameData.json}]
{
  "Id": "flame_strike",
  "Name": "Flame Strike",
  "Type": "Attack",
  "ManaCost": 2,
  "StarRating": 3,
  "Actions": [
    { "Type": "DealDamage", "TargetType": "AllEnemies", "Value": 8 },
    { "Type": "ApplyEffect", "EffectId": "burn", "Value": 3 }
  ]
}
\end{lstlisting}

\subsection{Heuristic AI Implementation}

The deterministic \texttt{HeuristicPlayerAI} implements \texttt{IPlayerAgent} with the decision logic described in Chapter~2. The key helper methods are:

\begin{lstlisting}[language=C, caption={Combat decision logic of HeuristicPlayerAI}]
public CombatDecision GetCombatDecision(GameRun run)
{
    var hero = run.TheHero;
    var enemies = run.CurrentCombat.Enemies
        .Where(e => e.CurrentHealth > 0).ToList();
    int incomingDamage = CalculateIncomingDamage(run, enemies);
    int neededBlock = Math.Max(0, incomingDamage - hero.Block);

    // Priority 1: Survival
    if (neededBlock > 0)
    {
        var def = FindBestDefensiveMove(run, hero, enemies, neededBlock);
        if (def.HasValue) return def.Value;
    }
    // Priority 2: Lethal
    var lethal = FindLethalMove(run, hero, enemies);
    if (lethal.HasValue) return lethal.Value;
    // Priority 3: Best general move
    return GetBestGeneralMove(run, hero, enemies);
}
\end{lstlisting}

\section{Cross-Platform C\# -- Python Integration}

\subsection{Motivation and Architecture}

The initial integration approach used gRPC with an Ngrok tunnel to forward requests from Kaggle (Linux cloud) to a local Windows server running the game engine. Benchmarking revealed a latency of $\sim$25\,ms per call, making the target of 1 million training steps infeasible (it would require $\sim$7 hours of network time alone).

The solution is \textbf{pythonnet} -- a library that embeds the .NET CLR directly inside the Python process, enabling zero-copy, function-call-speed interoperability. The key challenge was building the C\# project as a cross-platform .NET 9 DLL that can run on Linux and loading it inside a Kaggle notebook.

\subsection{Kaggle Setup Script}

The following setup is executed at the beginning of each Kaggle notebook:

\begin{lstlisting}[language=Python, caption={Kaggle environment setup for pythonnet}]
import subprocess, os

# Install .NET 9 SDK on Kaggle Linux
subprocess.run(["apt-get", "install", "-y", "dotnet-sdk-9.0"],
               capture_output=True)

# Set environment variables for pythonnet CoreCLR
os.environ["DOTNET_ROOT"] = "/usr/share/dotnet"
os.environ["PYTHONNET_PYDLL"] = "/usr/lib/x86_64-linux-gnu/libpython3.10.so"

import pythonnet
pythonnet.load("coreclr")

import clr
clr.AddReference("/kaggle/working/Roguelike.Core.dll")
from Roguelike.Core import GameController, GameDataLoader
\end{lstlisting}

After setup, calling C\# methods from Python incurs only the cost of CLR managed-to-native marshaling, typically $<1\,\mu$s per call -- a $25,000\times$ improvement over the gRPC approach.

\section{Reinforcement Learning Environment}

\subsection{DotNetGameEnv (Gymnasium Wrapper)}

The \texttt{DotNetGameEnv} class in \texttt{python\_roguelike/env/roguelike\_env.py} subclasses \texttt{gymnasium.Env} and wraps the C\# \texttt{GameController}. Its key methods are:

\begin{lstlisting}[language=Python, caption={Core Gymnasium environment methods}]
OBS_DIM = 148
ACT_DIM = 68

class DotNetGameEnv(gym.Env):
    def reset(self, seed=None, options=None):
        self._controller = GameController(self._load_game_data())
        obs = self._get_obs()
        return obs, {}

    def step(self, action: int):
        # Execute action in C# game engine
        result = self._controller.TakeAction(action)
        obs = self._get_obs()
        reward = self._compute_reward(result)
        done = result.IsTerminal
        info = {"win": result.IsVictory, "floor": result.Floor}
        return obs, reward, done, False, info

    def action_masks(self) -> np.ndarray:
        """Return binary mask of valid actions (required by MaskablePPO)."""
        valid = self._controller.GetValidActions()
        mask = np.zeros(ACT_DIM, dtype=bool)
        for a in valid:
            mask[a] = True
        return mask
\end{lstlisting}

\subsection{Domain Randomization}

The \texttt{reset()} method applies DR when enabled:

\begin{lstlisting}[language=Python, caption={Domain Randomization in environment reset}]
def _load_game_data(self):
    base = self._base_game_data  # loaded once at init
    if self._dr_noise <= 0.0:
        return base
    # Apply uniform noise to each numerical parameter
    noisy = copy.deepcopy(base)
    noise = self._dr_noise
    for card in noisy.Cards:
        for action in card.Actions:
            if action.Value > 0:
                action.Value = round(
                    action.Value * np.random.uniform(1-noise, 1+noise))
    for enemy in noisy.Enemies:
        enemy.MaxHealth = max(1, round(
            enemy.MaxHealth * np.random.uniform(1-noise, 1+noise)))
        for action in enemy.Actions:
            if action.Value > 0:
                action.Value = round(
                    action.Value * np.random.uniform(1-noise, 1+noise))
    return noisy

def set_dr_noise(self, noise: float):
    """Externally control DR level (used by curriculum callback)."""
    self._dr_noise = max(0.0, min(noise, 0.20))
\end{lstlisting}

A \texttt{DRCurriculumCallback} (in \texttt{training\_utils.py}) ramps DR noise from 0\% to 8\% linearly over the first 500,000 training steps of Phase~2, preventing the performance shock of instantly switching to high noise.

\subsection{Observation Construction}

The \texttt{\_get\_obs()} method queries the C\# game state and serializes it into the 148-dimensional NumPy array. A simplified excerpt:

\begin{lstlisting}[language=Python, caption={Observation vector construction}]
def _get_obs(self) -> np.ndarray:
    obs = np.zeros(OBS_DIM, dtype=np.float32)
    run = self._controller.CurrentRun
    hero = run.TheHero

    # Hero state (indices 0-3)
    obs[0] = hero.CurrentHealth / hero.MaxHealth
    obs[1] = hero.Block / max(hero.MaxHealth, 1)
    obs[2] = hero.Mana / max(hero.MaxMana, 1)
    obs[3] = hero.Gold / 200.0  # normalized by max expected gold

    # Active effects (indices 4-10)
    for i, eff in enumerate(EFFECT_LIST[:7]):
        obs[4 + i] = run.GetEffectIntensity(eff) / 10.0

    # Hand cards (indices 11-80)
    for slot, card in enumerate(run.Hand[:10]):
        base = 11 + slot * 7
        obs[base]   = card.Type / 3.0
        obs[base+1] = card.StarRating / 5.0
        obs[base+2] = card.ManaCost / 5.0
        obs[base+3] = card.PrimaryValue / 30.0
        obs[base+4] = float(card.CanTarget)
        obs[base+5] = float(card.Exhaust)
        obs[base+6] = float(hero.Mana >= card.ManaCost)
    # ... (enemies, relics, map, new features)
    return obs
\end{lstlisting}

\section{Agent Training Pipeline}

\subsection{Two-Phase Training Strategy}

Each agent is trained in two phases:
\begin{enumerate}
    \item \textbf{Phase 1 -- Stable Learning (4 million steps)}: DR is disabled ($\delta = 0$). The agent learns on the nominal parameter set, establishing a reliable policy from scratch. LR starts at $3\times10^{-4}$ and follows a linear schedule down to $1\times10^{-5}$.
    \item \textbf{Phase 2 -- Robustness Training (4 million steps, total 8 million)}: DR is enabled with a curriculum ramp from $\delta=0$ to $\delta=0.08$ over the first 500K steps. LR uses a cosine decay from $3\times10^{-5}$ to $5\times10^{-6}$.
\end{enumerate}
The Phase 2 anti-oscillation settings (reduced epochs 6$\to$4, tighter clip 0.2$\to$0.15, boosted entropy coefficient 0.05$\to$0.08) were determined empirically from analyzing the training logs where performance degraded after the abrupt DR introduction in earlier training runs.

\subsection{MaskablePPO Hyperparameters}

\begin{table}[H]
    \centering
    \caption{MaskablePPO hyperparameters for all four agents.}
    \label{tab:ppo_params}
    \begin{tabular}{lcc}
        \toprule
        \textbf{Parameter} & \textbf{Phase 1} & \textbf{Phase 2} \\
        \midrule
        Network architecture  & MLP [512, 512] & MLP [512, 512] \\
        Learning rate         & $3\times10^{-4}$ (linear $\to 1\times10^{-5}$) & Cosine $3\times10^{-5} \to 5\times10^{-6}$ \\
        n\_steps (rollout)    & 2048  & 2048 \\
        Batch size            & 256   & 256 \\
        n\_epochs             & 10    & 4 \\
        Clip range ($\epsilon$) & 0.20 & 0.15 \\
        Entropy coefficient   & 0.05  & 0.08 \\
        GAE $\lambda$         & 0.95  & 0.95 \\
        Discount $\gamma$     & 0.99  & 0.99 \\
        Parallel environments & 4 (SubprocVecEnv) & 4 \\
        Training steps        & 4M    & 4M \\
        DR noise $\delta$     & 0.00  & 0.00 $\to$ 0.08 \\
        \bottomrule
    \end{tabular}
\end{table}

\subsection{Deck Quality Reward}

All four agents include a shared \textbf{deck quality reward} at card selection events. This addresses the problem where agents selecting weak early-game cards (1--2 stars) fail to build a deck capable of defeating late-game content:
\begin{equation}
r_{\text{deck}} = \begin{cases}
+8 & \text{if selected card is 4 or 5 stars} \\
+3 & \text{if selected card is 3 stars} \\
0  & \text{otherwise (1 or 2 stars)}
\end{cases}
\end{equation}
Cards of 1 or 2 stars receive no reward (rather than a penalty) because they are necessary and valid choices early in a run.

\subsection{Training Infrastructure}

Training was executed on Kaggle's free GPU tier (NVIDIA Tesla T4, 16 GB VRAM) using 4 parallel environments (\texttt{SubprocVecEnv}). Each agent's 8-million-step training run required approximately 6--8 hours. Models were checkpointed every 500K steps using \texttt{MaskableEvalCB}, which evaluates on 20 deterministic episodes (no DR noise, fixed seeds) and saves the best checkpoint by win rate.

\begin{lstlisting}[language=Python, caption={Training launch code (simplified)}]
from sb3_contrib import MaskablePPO
from python_roguelike.training.training_utils import (
    DRCurriculumCallback, MaskableEvalCB, cosine_lr_fn
)

# Phase 1
model = MaskablePPO(
    "MlpPolicy", env,
    learning_rate=3e-4,
    n_steps=2048, batch_size=256, n_epochs=10,
    ent_coef=0.05, clip_range=0.20,
    policy_kwargs=dict(net_arch=[512, 512]),
    verbose=1
)
model.learn(total_timesteps=4_000_000,
            callback=MaskableEvalCB(eval_env, n_eval_episodes=20))
model.save("phase1_checkpoint")

# Phase 2 -- load Phase 1 weights
model = MaskablePPO.load("phase1_checkpoint", env=env)
model.learning_rate = cosine_lr_fn(3e-5, 5e-6)
model.ent_coef = 0.08
model.clip_range = 0.15
model.n_epochs = 4
dr_callback = DRCurriculumCallback(
    env, target_noise=0.08, ramp_steps=500_000)
model.learn(total_timesteps=4_000_000,
            callback=[dr_callback,
                      MaskableEvalCB(eval_env, n_eval_episodes=20)])
model.save("final_checkpoint")
\end{lstlisting}

\section{Evolutionary Optimization Implementation}

\subsection{Single-Objective Genetic Algorithm}

The GA is implemented in C\# (\texttt{GeneticAlgorithm.cs}) for performance. The main evolutionary loop:
\begin{enumerate}
    \item Initialize population of 100 \texttt{BalanceGenome} objects with all scalars drawn from $\mathcal{U}(0.7, 1.3)$.
    \item Evaluate fitness: for each genome, run 200 complete game simulations using the DHA; aggregate win rate, victory HP, and floor-on-death; compute $F(\mathbf{g})$ via Equation~\ref{eq:ga_fitness} with weights $w_1=0.5, w_2=0.3, w_3=0.2$.
    \item Tournament selection (size $k=3$) to fill the mating pool.
    \item Uniform crossover ($p_c = 0.80$) to produce offspring.
    \item Gaussian mutation ($p_m = 0.02$, $\sigma = 0.10$).
    \item Elitism: copy top 10\% of previous population unchanged.
    \item Repeat for 30 generations.
\end{enumerate}
Evaluation is parallelized using \texttt{Parallel.ForEach} over all 8 logical cores, reducing wall-clock time from $\sim$9 hours to $\sim$73 minutes.

\subsection{NSGA-II Multi-Objective Optimizer}

The NSGA-II is implemented in \texttt{NSGA2Optimizer.cs} following Deb et al.\ \cite{deb2002fast}:

\begin{lstlisting}[language=C, caption={Fast non-dominated sorting (core NSGA-II operation)}]
public List<List<Individual>> FastNonDominatedSort(
    List<Individual> population)
{
    var domCount = population.ToDictionary(p => p, _ => 0);
    var dominated = population.ToDictionary(p => p,
                                _ => new List<Individual>());
    var firstFront = new List<Individual>();

    foreach (var p in population)
        foreach (var q in population)
        {
            if (p == q) continue;
            if (p.Fitness.Dominates(q.Fitness))
                dominated[p].Add(q);
            else if (q.Fitness.Dominates(p.Fitness))
                domCount[p]++;
        }

    foreach (var p in population)
        if (domCount[p] == 0) { p.Fitness.Rank = 1; firstFront.Add(p); }

    var fronts = new List<List<Individual>> { firstFront };
    int rank = 1;
    while (fronts[rank - 1].Any())
    {
        var next = new List<Individual>();
        foreach (var p in fronts[rank - 1])
            foreach (var q in dominated[p])
                if (--domCount[q] == 0)
                { q.Fitness.Rank = rank + 1; next.Add(q); }
        if (next.Any()) { fronts.Add(next); rank++; }
        else break;
    }
    return fronts;
}
\end{lstlisting}

The \texttt{HierarchicalGenome} uses SBX crossover with distribution index $\eta_c = 20$ and polynomial mutation with $\eta_m = 20$, as recommended for continuous-variable optimization \cite{deb2002fast}. The adaptive evaluation strategy runs 50 simulations for initial fitness estimation and up to 200 for genomes that appear on or near the Pareto front, balancing accuracy with computational cost.

\subsection{Python-Side Fitness Evaluation for RL Agents}

When RL agents are used as evaluators (post-training), the fitness evaluation moves to Python:

\begin{lstlisting}[language=Python, caption={RL-agent-based fitness evaluation}]
def evaluate_genome(genome, agents, n_runs=50):
    """Evaluate a genome using the 4 trained RL agents."""
    game_data = apply_genome(BASE_GAME_DATA, genome)
    stats = {name: [] for name in agents}

    for name, agent in agents.items():
        env = DotNetGameEnv(game_data, dr_noise=0.0)
        for seed in range(n_runs):
            obs, _ = env.reset(seed=seed)
            done = False
            while not done:
                action, _ = agent.predict(
                    obs, action_masks=env.action_masks(),
                    deterministic=True)
                obs, _, done, _, info = env.step(action)
            stats[name].append(info)

    return compute_fitness(stats)
\end{lstlisting}

This RL-based evaluation is used for the final validation runs and the card viability analysis in Chapter~4.
